% ============================================================================
% Fichier  : partie3/P03_C10_iso27002.tex
% Titre    : ISO/IEC 27002 et les Mesures de Sécurité
% Auteur   : MINKA MI NGUIDJOÏ Thierry Emmanuel
% ============================================================================
% Sources : ISO/IEC 27002:2022 ; ISO/IEC 27001:2022, Annexe A. Fictif : Eséka.
% ============================================================================

\chapter{ISO/IEC 27002 et les Mesures de Sécurité}
\label{chap:iso27002}
\chaptermark{ISO 27002 et l'Annexe A}

\epigraph{%
    « Les mesures ne se choisissent pas dans un catalogue par goût~:\\
    elles se justifient par le risque. »%
}{Principe du cours}

\niveauChapitre{Intermédiaire}{Le Chapitre~\ref{chap:iso27000}.}

\begin{boxobjectifs}
\begin{itemize}
    \item Situer \norme{ISO/IEC 27002} et l'Annexe~A de la 27001.
    \item Énoncer les quatre thèmes de mesures de la version 2022.
    \item Comprendre la structure d'une mesure (attributs, objectif).
    \item Relier mesures et appréciation des risques.
    \item Introduire la déclaration d'applicabilité.
\end{itemize}
\end{boxobjectifs}

\bigskip

% ============================================================================
\section*{Question directrice : parmi toutes les mesures possibles, lesquelles retenir ?}
% ============================================================================

La 27002 décrit des dizaines de mesures. La mairie ne les appliquera pas
toutes~: elle retiendra celles que son risque justifie. Ce chapitre présente le
catalogue et prépare le choix.

\medskip
\textbf{Plan du chapitre.}
\begin{enumerate}
    \item 27002 et Annexe~A (section~\ref{sec:c10-annexea}).
    \item Les quatre thèmes (section~\ref{sec:c10-themes}).
    \item Structure d'une mesure (section~\ref{sec:c10-structure}).
    \item Mesures et risques (section~\ref{sec:c10-risques}).
    \item Vers la déclaration d'applicabilité (section~\ref{sec:c10-soa}).
\end{enumerate}

\separateur

% ============================================================================
\section{27002 et Annexe A}
\label{sec:c10-annexea}
% ============================================================================

\begin{boxdef}
L'\termecle{Annexe~A} de \norme{ISO/IEC 27001} liste les mesures de sécurité de
référence. La norme \norme{ISO/IEC 27002} en donne le \termecle{guide de mise
en \oe uvre}~: pour chaque mesure, son objectif et des conseils d'application.
\end{boxdef}

L'Annexe~A fournit la liste de contrôle~; la 27002 explique comment faire. La
version 2022 réorganise les mesures en quatre thèmes et introduit des
\termecle{attributs} (type de mesure, propriété de sécurité, fonction).

% ============================================================================
\section{Les Quatre Thèmes}
\label{sec:c10-themes}
% ============================================================================

\begin{tableaucours}{L{3.4cm}L{2.0cm}L{6.6cm}}{%
    \tableauHeader{Thème} &
    \tableauHeader{Nombre} &
    \tableauHeader{Exemples}
}
Organisationnelles & 37 & Politiques, rôles, fournisseurs, incidents \\
Liées aux personnes & 8 & Sensibilisation, filtrage, télétravail \\
Physiques & 14 & Accès aux locaux, supports, équipements \\
Technologiques & 34 & Accès, chiffrement, journalisation, sauvegarde \\
\end{tableaucours}
\vspace{-0.8em}
\begin{center}{\small\itshape Tableau~10.1. Les quatre thèmes de l'ISO/IEC 27002:2022}\end{center}

Au total, 93~mesures. Les parties V du cours détaillent les mesures
organisationnelles et humaines (Chapitre~\ref{chap:orga}), techniques
(Chapitre~\ref{chap:technique}) et de continuité (Chapitre~\ref{chap:continuite}).

% ============================================================================
\section{Structure d'une Mesure}
\label{sec:c10-structure}
% ============================================================================

\begin{boxprocedure}
\textbf{Constatation.} La mesure « sauvegarde de l'information » (27002) a pour
objectif de garantir la \emph{disponibilité}~; ses attributs la classent en
mesure préventive et corrective.

\medskip
\textbf{Analyse.} Chaque mesure a un objectif (ce qu'elle protège) et un guide
(comment l'appliquer). À la mairie, cette mesure impose une règle de
sauvegarde quotidienne et un test de restauration (Chapitre~\ref{chap:continuite}).
\end{boxprocedure}

Comprendre la structure évite l'application mécanique~: on applique une mesure
pour son \emph{objectif}, pas pour cocher une case.

% ============================================================================
\section{Mesures et Risques}
\label{sec:c10-risques}
% ============================================================================

La norme est claire~: on ne retient pas une mesure parce qu'elle figure dans
l'Annexe~A, mais parce que l'\termecle{appréciation des risques} (partie IV) la
justifie. Les mesures sont la réponse au risque, pas une liste à appliquer en
bloc.

\begin{boxmetier}{Auditeur SMSI}
\textbf{La traçabilité risque-mesure.} L'auditeur vérifie que chaque mesure
retenue répond à un risque identifié, et que chaque risque significatif reçoit
une mesure. Une mesure sans risque et un risque sans mesure sont deux écarts.
\end{boxmetier}

% ============================================================================
\section{Vers la Déclaration d'Applicabilité}
\label{sec:c10-soa}
% ============================================================================

\begin{boxdef}
La \termecle{déclaration d'applicabilité} (\sigle{DdA}, \emph{Statement of
Applicability})~: document qui, pour chaque mesure de l'Annexe~A, indique si
elle est retenue, sa justification, et son état de mise en \oe uvre.
\end{boxdef}

La \sigle{DdA} est la pièce maîtresse qui relie les risques aux mesures. Elle
est détaillée au Chapitre~\ref{chap:dda}, après l'étude des risques. Retenir
ou écarter une mesure s'y justifie explicitement.

\separateur

\begin{boxresume}
\textbf{Ce qu'il faut retenir de ce chapitre~:}
\begin{itemize}
    \item L'\textbf{Annexe~A} liste les mesures~; la \textbf{27002} guide leur
          mise en \oe uvre.
    \item La version 2022 compte \textbf{93 mesures} en quatre thèmes.
    \item Chaque mesure a un \textbf{objectif}~; on l'applique pour cet
          objectif.
    \item Les mesures se justifient par le \textbf{risque}, non par le
          catalogue.
    \item La \textbf{déclaration d'applicabilité} relie risques et mesures.
    \end{itemize}
\end{boxresume}

\begin{boxinfo}
\textbf{Points de contrôle de l'élève.} Je sais~: distinguer Annexe~A et
27002~; citer les quatre thèmes~; décrire la structure d'une mesure~; relier
mesure et risque~; définir la \sigle{DdA}.
\end{boxinfo}

\bigskip

\begin{boxexo}[title={\ding{45}\quad Exercice~10.1. Annexe A ou 27002 \niveaubase}]
Laquelle liste les mesures, laquelle en guide la mise en \oe uvre~?
\end{boxexo}

\begin{boxexo}[title={\ding{45}\quad Exercice~10.2. Les quatre thèmes \niveaubase}]
Citez les quatre thèmes de mesures et un exemple pour chacun.
\end{boxexo}

\begin{boxexo}[title={\ding{45}\quad Exercice~10.3. Objectif d'une mesure \niveaubase}]
Quel est l'objectif de la mesure « sauvegarde de l'information »~?
\end{boxexo}

\begin{boxexo}[title={\ding{45}\quad Exercice~10.4. Mesure et risque \niveaumoyen}]
Pourquoi n'applique-t-on pas toutes les mesures de l'Annexe~A~?
\end{boxexo}

\begin{boxexo}[title={\ding{45}\quad Exercice~10.5. Classer des mesures \niveaumoyen}]
Classez par thème~: chiffrement, sensibilisation, contrôle d'accès aux locaux,
gestion des fournisseurs.
\end{boxexo}

\begin{boxexo}[title={\ding{45}\quad Exercice~10.6. Deux écarts \niveaumoyen}]
Qu'est-ce qu'une « mesure sans risque » et un « risque sans mesure »~?
\end{boxexo}

\begin{boxexo}[title={\ding{45}\quad Exercice~10.7. Mesures pour la mairie \niveaumoyen}]
Citez trois mesures de l'Annexe~A particulièrement utiles à la mairie d'Eséka.
\end{boxexo}

\begin{boxexo}[title={\ding{45}\quad Exercice~10.8. Rôle de la DdA \niveaumoyen}]
À quoi sert la déclaration d'applicabilité~?
\end{boxexo}

\begin{boxexo}[title={\ding{45}\quad Exercice~10.9. Fiche de mesure \niveauexpert}]
Rédigez une fiche de mesure pour « contrôle d'accès » (objectif, mise en
\oe uvre à la mairie, preuve attendue).
\end{boxexo}

\begin{boxexo}[title={\ding{45}\quad Exercice~10.10. Extrait de DdA \niveauexpert}]
Rédigez un extrait de \sigle{DdA} pour cinq mesures (retenue~? justification~?
état~?).
\end{boxexo}

\bigskip

\begin{boxinfo}
\textbf{Pour aller plus loin.}
\begin{itemize}
    \item ISO/IEC 27002~\cite{iso27002}~; ISO/IEC 27001, Annexe~A~\cite{iso27001}.
\end{itemize}
\end{boxinfo}

\begin{boxinfo}
\textbf{Transition.} Choisir des mesures suppose de connaître ses risques~: la
partie IV s'ouvre au Chapitre~\ref{chap:vocabrisque} sur le vocabulaire du
risque.
\end{boxinfo}
